\section{Calculations as requests}\label{sec:requests}

A record says what each statement rests on. It does not perform calculations. In a project carried out with language models, the
calculations are the part most likely to be invented: a model asked for a number will produce one. QAM therefore separates the request
for a calculation from its execution and from the stored result. Requests are written in a small request language, MML; they are
executed by an executor, MMA; and every execution leaves a \emph{receipt}.

\subsection{Three outcomes}

The central design decision is that an execution has three possible outcomes, not two:
\begin{itemize}[leftmargin=1.5em]
\item \cont: the request was admissible and was carried out; the receipt contains the result.
\item \guard: the request was well formed, but a condition it depends on is not available or not fulfilled, for example a premise that is not in the
record. No result is returned or stored. The receipt names the missing condition as an obligation.
\item \reject: the request itself is not admissible, for example because an argument lies outside the declared range.
\end{itemize}
A \guard\ is not a failure. It is the third outcome required by the decidability boundary of Section~\ref{sec:intro}: the honest answer
when a precondition is open. It is stored, hashed and retrievable like any other result.

\subsection{The example as a request}

The interval calculation of Section~\ref{sec:example} is available as an allowlisted recipe. A request names the recipe, the frames, the
preparation, the boost and the measured interval:
\begin{lstlisting}
{ "id": "gram_boost_interval", "instruction": "READOUT", "guards": [],
  "arguments": { "recipe": "em.gram-boost-range",
                 "preparation": "E2x-B-unit-yz", "frame": "Lab",
                 "target_frame": "Boosted", "boost": "z:1/4",
                 "a": "-1/4", "b": "1/4" } }
\end{lstlisting}
Before computing, the executor looks up the analytic rule the recipe depends on, by the hash of its record entry, and follows the
premises of that entry recursively. We ran the request four times, the first as shown and then with one thing changed each time:
\begin{center}
\small
\begin{tabular}{@{}p{0.34\textwidth}p{0.12\textwidth}p{0.45\textwidth}@{}}
\toprule
Request & Outcome & Receipt states\\
\midrule
as above, rule indexed & \cont & target interval $[121/162,\ 961/1202]$\\
rule not in the database & \guard & obligation: missing premise (the rule's hash)\\
interval with $a>b$ & \reject & error: endpoints must satisfy $-1\le a\le b\le 1$\\
boost $z{:}1/3$ instead of $z{:}1/4$ & \reject & error: the recipe requires the declared preparation and boost\\
\bottomrule
\end{tabular}
\end{center}
The first row reproduces the hand calculation~\eqref{eq:interval} with exact rational arithmetic. The second shows the guard: the same
request against a database in which the rule is not recorded returns no result, because the result would rest on something the record does
not contain. The third and fourth show that the recipe accepts only the situation for which the rule was derived; it does not
extrapolate to another boost speed.

\subsection{What a receipt contains}

A receipt is a structured document identified by its SHA-512 hash. It contains the hash of the request (the request itself is stored beside it); the outcome; the result or the reasons
for blocking it; a hash of the context in which the request was evaluated; hashes of the implementation that executed it; and, when the
request depends on record entries, a \emph{premise trace}: a snapshot of every entry reached from the named premises, in dependency
order, with its statement, status and frame. For the example, the trace contains two entries, the preparation in Lab and the rule obtained
from it by conversion into Boosted, both \st{conditional}.

Two fields deserve attention. The receipt for the successful run sets
\begin{lstlisting}
"closure_proved": false,
"mathematical_relations_proved": false,
\end{lstlisting}
A successful calculation is not a proof that the rule it applies is correct. The rule has the status recorded for it, here
\st{conditional}, and the receipt does not change that. In the same spirit, a premise linked with status \st{reading} stays a reading
even when every calculation built on it succeeds.

A stored receipt can be checked in two ways. Verification recomputes the stored hashes. \emph{Replay} reruns the request with the stored
context and the current implementation, bypassing any cache, and reports separately whether the receipt and the result agree. For the
example, replay reports that both match. A changed implementation can change the receipt hash while the result stays the same; the
report keeps the two apart.

\subsection{The abstract machine}

The three-way execution has an abstract model that is small enough to state completely. It is the store machine defined in the main
text of the research project~\cite{repo}, with a small instruction set of its own (store, link, snapshot, rollback, scrub, export). The
running executor follows the same three-way pattern with a different instruction set; the \texttt{READOUT} instruction of the example
is not part of the model. A \emph{store} consists of its content (objects at
addresses and relations between addresses), a register of guards (predicates that the content must satisfy), a list of snapshots and a
log. A store is \emph{admissible} if every relation uses stored addresses and every guard holds. Each instruction $i$ has a raw effect,
which may be undefined, and is executed as
\[
\mathrm{exec}(s,i)=\begin{cases}
(s \text{ with log entry } (i,\reject),\ \reject) & \text{if the raw effect is undefined},\\
(s' \text{ with log entry } (i,\cont),\ \cont) & \text{if the raw effect } s' \text{ is admissible},\\
(s \text{ with log entry } (i,\guard),\ \guard) & \text{otherwise.}
\end{cases}
\]
A result that would violate a guard is never adopted; the old content stays, and the log records what happened. In the model, \guard\
thus means that a computed result was discarded. The running executor answers \guard\ also when a premise or a declared guard is
missing.

\begin{proposition}[Machine properties]\label{prop:machine}
For every store $s$ and instruction $i$:
\begin{enumerate}[label=(\roman*),nosep,leftmargin=2em]
\item if $s$ is admissible, so is the store after $\mathrm{exec}(s,i)$, and hence after any sequence of instructions;
\item the outcome is never a halt: every step ends in \cont, \guard\ or \reject;
\item if the outcome is not \cont, content, snapshots and guards are unchanged;
\item every step adds exactly one log entry, which records the instruction and its outcome.
\end{enumerate}
\end{proposition}

\begin{proof}
All four follow by inspecting the three cases of $\mathrm{exec}$. In the first and third case the content is that of $s$, which is
admissible by assumption; in the second it is $s'$, which is admissible by the case condition. The sequence statement follows by
induction. Every case prepends exactly one log entry and returns one of the three outcomes.
\end{proof}

A further rule concerns certificates. The executor can label a result as a \emph{candidate} for a theorem, but a certificate for a
statement $P$ is defined to contain a proof of $P$. No instruction of the model produces a proof, so a candidate becomes a certificate only through a proof supplied from outside the
machine. The Lean file records the immediate consequence of this definition, that there is no certificate for a contradiction; the
statement does not refer to the instructions (Appendix~\ref{app:lean}).

\paragraph{Implementation.} The running executor is a Python adapter over an SQLite database. It follows the same three-way
semantics and checks the declared guards of the context, the request and the objects involved, as well as the linked premises, before
any result is stored. That the adapter behaves like the abstract model is not proved, and its instruction set differs from the
model's. Its behaviour is supported by tests of the tape instruction, of batches and of receipts, and by the runs reported here. For one
instruction, the hardware comparison of Section~\ref{sec:tape} checks the adapter's results against other implementations,
which supports the results of that instruction but not the guard semantics as such.

\subsection{Batches}

For the finite permutation search used in one of the project's recipes, up to 64 requests can be run as a batch on one to four worker
processes. A request may be declared to depend on others; the dependency
graph is checked for cycles before the start. If a predecessor ends with \guard\ or \reject, or is itself blocked, its dependents are
not executed but listed as \emph{blocked}, with the predecessor that blocked them. A batch produces a hashed report in input order,
stored in the same database transaction as the individual receipts. Reports can be rerun, compared with earlier reports and audited for
consistent links without being rerun.
